\chapter{Literature Review}
\label{ch:lit}

An ANPR system chains four problems: finding the plate in the frame, restoring
a plate that the weather has degraded, reading its characters, and making a
decision from many frames rather than one. This chapter reviews the work that
the project builds on for each of them.

\section{Plate Detection}

Early plate locators used edge density, colour and morphology to find the
bright rectangle of a plate; they work on frontal, well-lit vehicles but fail
on tilted plates, clutter and bad weather. Modern systems treat the plate as an
object class for a single-stage detector. The YOLO family predicts boxes and
class scores in one pass over the image~\cite{redmon2016yolo}, which gives the
speed needed for video. YOLOv11~\cite{jocher2024yolo11} refines the backbone
with C3k2 blocks and is offered in sizes from the nano model upwards; the nano
variant is small enough for embedded hardware, and this is the detector used in
the project. Large plate datasets such as CCPD~\cite{xu2018ccpd} and the
RodoSol-ALPR set~\cite{laroca2022rodosol} show that detection of plates is
largely a solved problem in good conditions; the open questions are small
plates and degraded images.

\section{Plate Recognition}

General-purpose scene-text readers such as EasyOCR~\cite{jaided2020easyocr}
segment an image into text lines and recognise them with a recurrent network.
They are trained on signs and documents, not on plates, and in our experiments
they confuse plate characters (\plate{1}/\plate{I}, \plate{7}/\plate{Z}) and
fail on small, blurred plates. Plate-specific readers avoid explicit character
segmentation. LPRNet~\cite{zherzdev2018lprnet} is a compact convolutional
network that maps a plate crop to a sequence of per-column character
probabilities and is trained with the Connectionist Temporal Classification
(CTC) loss~\cite{graves2006ctc}, which sums over all alignments between the
column sequence and the target string. It reads the whole plate in one pass, is
small, and needs no recurrent layer, which makes it
suitable for real-time and embedded use. Its weakness, shared by all CTC
models, is that repeated characters need an intervening blank and are easily
shortened, and that it can only read what the training data contained.
Public plate datasets are dominated by Chinese and Brazilian plates, so
Indian-plate performance depends on self-collected or small public sets.

\section{Weather Restoration}

\subsection{Haze and fog}

Haze formation is described by the atmospheric scattering
model~\cite{narasimhan2002vision}
\begin{equation}
  I(x) = J(x)\,t(x) + A\,\bigl(1 - t(x)\bigr), \qquad t(x) = e^{-\beta d(x)},
  \label{eq:fog}
\end{equation}
where $I$ is the observed image, $J$ the clear scene, $t$ the transmission,
$A$ the atmospheric light, $\beta$ the scattering coefficient and $d$ the
distance of the point from the camera. Dehazing is the problem of recovering
$J$ from $I$ alone.

The \emph{Dark Channel Prior} of He \etal~\cite{he2011dcp} estimates $A$ and
$t$ from the observation that in haze-free outdoor patches at least one colour
channel is close to zero. It needs no training and works on light, uneven haze,
but its assumption fails on dense, uniform fog, where the whole frame, including
the subject, is veiled. \emph{CLAHE}~\cite{zuiderveld1994clahe} equalises the
histogram tile by tile with a clip limit; it contains no haze model, is very
fast, and can raise local contrast, but it cannot restore detail that was
never captured. The \emph{guided filter}~\cite{he2013guided} is an edge-preserving
smoother used here both to refine transmission maps and to clean network
output.

Learned dehazing networks are trained on pairs of hazy and clear images.
Large synthetic sets such as RESIDE~\cite{li2019reside} and Foggy
Cityscapes~\cite{sakaridis2018foggy} are generated from clear images with
Equation~\eqref{eq:fog} and a depth map, while NH-HAZE~\cite{ancuti2020nhhaze}
provides a small number of real, non-homogeneous fog pairs captured with a
fog machine. State space models, and in particular
Mamba~\cite{gu2023mamba}, process an image as a sequence with linear
complexity; UVM-Net~\cite{zheng2024uvmnet} places Mamba blocks in a U-shaped
encoder--decoder and reports strong dehazing at a moderate cost. UVM-Net is
the network (UVMNet) used in this project. Depth estimators such as Depth
Anything~V2~\cite{yang2024depthanythingv2} make it possible to generate
synthetic fog with realistic depth dependence from any clear photograph,
which we use in Section~\ref{sec:depthfog}.

A recurring finding in the literature and in our own work is the
\emph{synthetic-to-real gap}: networks trained on synthetic fog score high
PSNR on synthetic test sets but may leave real fog untouched or introduce
colour casts, and PSNR gains do not automatically translate into better
recognition by a downstream reader.

\subsection{Rain}

Rain adds streaks, spray and glare. Recent deraining networks use multi-scale
and implicit representations; NeRD-Rain~\cite{chen2024nerdrain}
exchanges information between coarse and fine scales and is used here as a
pretrained comparison baseline. NAFNet~\cite{chen2022nafnet} showed that a
simple baseline without nonlinear activations can match far more complex
restoration networks and is a design reference for lightweight models. As with
fog, deraining is trained on synthetic rain, and real rain also contains road
splash, wet-surface sheen and lens drops that synthetic streaks do not model.

\section{Multi-Frame Evidence}

A moving vehicle appears in many frames, each a slightly different view of the
same plate. Reading the plate in every frame and taking a vote, or aligning
the plate crops and combining them, uses this redundancy. Image alignment with
the enhanced correlation coefficient criterion~\cite{evangelidis2008ecc} gives
sub-pixel registration of crops under affine motion, after which the crops can
be averaged to reduce noise. Tracking plates across frames also lets the system
announce each vehicle once instead of once per frame.

\section{Summary}

The literature supports the design adopted here: a small single-stage detector,
a CTC-based plate reader trained on Indian plates, a learned dehazing network
for fog and a lightweight network for rain, and voting over frames for
reliability. It also identifies the open problems that Phase~1 addresses:
dehazing that helps recognition and not only appearance, readers that
generalise to Indian plate layouts, and a real-time implementation.
